\RequirePackage{fix-cm}
\documentclass[11pt,a4paper]{ltjsarticle}
\usepackage[margin=24mm]{geometry}
\usepackage{amsmath,amssymb,amsthm,mathtools,booktabs,longtable,array}
\usepackage{tikz-cd}
\usepackage[unicode,hidelinks]{hyperref}
\usepackage{bookmark}
\setlength{\emergencystretch}{3em}
\setlength{\parskip}{0.22em}
\newtheorem{theorem}{定理}[section]
\newtheorem{proposition}[theorem]{命題}
\newtheorem{lemma}[theorem]{補題}
\newtheorem{corollary}[theorem]{系}
\theoremstyle{definition}
\newtheorem{definition}[theorem]{定義}
\newtheorem{assumption}[theorem]{仮定}
\newtheorem{remark}[theorem]{注意}
\newtheorem{example}[theorem]{例}
\renewcommand{\proofname}{証明}
\newcommand{\cA}{\mathcal A}\newcommand{\cB}{\mathcal B}
\newcommand{\cC}{\mathcal C}\newcommand{\cD}{\mathcal D}\newcommand{\cH}{\mathcal H}
\newcommand{\cT}{\mathcal T}\newcommand{\cF}{\mathcal F}
\newcommand{\cZ}{\mathcal Z}\newcommand{\OO}{\mathcal O}
\newcommand{\Db}{D^b}\newcommand{\Sfun}{\mathbb S}
\newcommand{\AR}{\mathcal A_{\mathrm R}}
\newcommand{\AV}{\mathcal A_{\mathrm V}}
\newcommand{\AAS}{\mathcal A_{\mathrm{AS}}}
\newcommand{\AAC}{\mathcal A_{\mathrm{ac}}}
\DeclareMathOperator{\Hom}{Hom}\DeclareMathOperator{\Ext}{Ext}
\DeclareMathOperator{\End}{End}\DeclareMathOperator{\RHom}{RHom}
\DeclareMathOperator{\coh}{coh}\DeclareMathOperator{\qgr}{qgr}
\DeclareMathOperator{\QGr}{QGr}\DeclareMathOperator{\gr}{gr}
\DeclareMathOperator{\Perf}{Perf}\DeclareMathOperator{\Pic}{Pic}
\DeclareMathOperator{\Ind}{Ind}\DeclareMathOperator{\thick}{thick}
\DeclareMathOperator{\rk}{rk}\DeclareMathOperator{\add}{add}
\DeclareMathOperator{\Ann}{Ann}
\title{非可換del Pezzo曲面の四つのheart\\
\large VdB・Rains・反標準・Artin--Schelter構成の比較}
\author{}
\date{2026年11月16日\qquad 統合研究稿 hearts-v17}
\hypersetup{pdftitle={非可換del Pezzo曲面の四つのheart},pdfsubject={導来比較、HRS tilt、共鳴、反標準corner}}
\begin{document}
\maketitle
\begin{abstract}
非可換del Pezzo曲面に付随するVdB heart，Rains heart，anticanonical heart，
AS heartを，比較関手と構造対象を指定して比較する．同じblow-up表示に由来する
VdBとRainsのt構造は，blow-down関手のpseudo-canonical性によって同定される．
一方，blow-upの各段階で行うsection modelへの移行は別の操作であり，
Van den Berghの適用条件下では一回のHRS tiltとして記述される．
滑らかなgenuine del Pezzo曲面と滑らかな楕円境界からのRains gluingについては，
自然な反標準導来切断を用いて導来同値と全heartの振幅$[0,1]$を示す．
負指数の実在共鳴根は，導来同値を失わずに反標準AZ amplenessとheart一致が失敗することを検出する．
AS側では，helix全体の豊富性と一つの構造対象の反Serre orbitの豊富性を区別し，
後者をnoetherianなfolded algebraのidempotent cornerによって判定する．
三次曲面の混合$A_2$部分系を例として，connecting map，有限長部分圏，
反標準側の三角表現と四次元アフィン商を記述する．
一般の四heartの無条件な一致，AS比較の全域的t-exactness，および族での降下は主張しない．
\end{abstract}
\tableofcontents

\section{序論：四つの構成と三つの比較問題}\label{sec:intro}
滑らかな可換del Pezzo曲面では，点のblow-up，反標準環，例外列から作るhelixは
いずれも同じ圏$\coh S$を回復する．非可換化すると，この一致には二つの障害が現れる．
第一に，同じ導来圏の中で標準的なheartが変わり得る．第二に，導来同値が存在しても，
構造対象$\OO$をどの対象へ送るかが自動的には決まらない．
本稿では次の三つを分ける：abel圏の構成，指定された導来同値，その同値のt-exactnessである．
構造対象・反Serre作用・境界restriction・標識の保存はさらに別に確認する．

名称は構成の由来を表す．\emph{VdB heart}はVan den Berghの実際のblow-upで得る
coherent sheaf圏，\emph{Rains heart}はRainsのgluingとblowdown表示で指定するt構造のheart，
\emph{anticanonical heart}は指定対象の反標準環の$\qgr$，
\emph{AS heart}はArtin--Schelter正則$\mathbb Z$-代数の$\qgr$である．
四名称が四つの相異なる圏を意味するわけではない．特に同じblow-up表示を用いる場合の
VdB--Rains一致と，VdBのsection modelとの比較は混同してはならない．

\subsection{主結果の見取り図}
本稿の比較は次の形を取る．各矢の仮定は後の節で明記する．
\begin{equation}\label{eq:overview}
\begin{tikzcd}[column sep=large,row sep=large]
\AV \arrow[r,"\text{同じblow-up表示}","\simeq"'] & \AR
\arrow[d,phantom,"\text{HRS tilt}" description] \\
\AAS \arrow[r,"\text{有限corner条件}","\simeq"'] & \AAC
\end{tikzcd}
\end{equation}
右側の縦の表示はabel同値の矢ではない．指定した導来同値で同じ三角圏に入れたとき，
一方のheartが他方のHRS tiltになることを表す．下側の横矢は条件付きであり，
任意のAS表示に存在する同値ではない．

\begin{theorem}[比較定理の要約]\label{thm:summary}
以下の各主張を，それぞれの適用範囲で得る．
\begin{enumerate}
\item 同じ初期heart，blow-up中心，gluing bimodule，blow-down関手を同定した
VdB表示とRains表示では，誘導される導来同値はt-exactである
（定理\ref{thm:RV}）．
\item 滑らかなgenuine可換del Pezzo $S$，滑らかな$E\in|-K_S|$，
$q\in\Pic^0(E)$からのRains gluingでは，自然な関手
$\Phi:\Db(\AR)\simeq\Db(\AAC)$がある．$H=-K$がnefなRains heartに対し
\[
\Phi^{-1}(\AAC)=\langle\cF[1],\cT\rangle,
\quad \cT=\ker H^1\Phi,\quad\cF=\ker H^0\Phi
\]
である（定理\ref{thm:derived}，\ref{thm:globaltilt}）．
\item 境界非交叉の既約根sheaf $M$が実在し，$m=\chi(\OO,M)<0$なら，
$(\OO,\Sfun^{-1}[2])$はRains heartでAZ-ampleでない．自然な$\Phi$はt-exactでないが，
(2)の範囲では依然として導来同値である（系\ref{cor:negative}）．
\item AS helixをfoldした代数$B$と構造頂点$p$が
命題\ref{prop:corner}の有限性条件を満たし，$pBp$が指定した反標準環であるとする．
自然な$\AAS\to\AAC$がabel同値となる必要十分条件は
$\dim_k B/BpB<\infty$である（系\ref{cor:ASac}）．
\end{enumerate}
\end{theorem}
さらに三次曲面型の標準noetherian周期表示では，(4)の有限性条件と指定AOU同値のt-exactnessを
四交比の十一の非零条件で判定する（定理\ref{thm:AOUexactlocus}）．
係数トーラス上の一致開集合は，古典的六点モジュライと四次元トーラスの積となる
（系\ref{cor:coefficientopen}）．これは係数空間の記述であり，全点のAS正則性を主張しない．

二つの追加的な接続を重視する．第一に，VdBの最後のsection-model関手は，
構造対象と反Serre作用を保つことから内在的な反標準導来切断と自然に同定される
（定理\ref{thm:VdBcanonical}）．従って同じ最終blow-upのRains実現では，
比較関手を別の仮定として与えずに四者を接続できる（系\ref{cor:fourfunctors}）．
第二に，境界非交叉部分のtorsion pairは全ての部分対象・商のEuler標数で決まり，
有限ADE型では負指数の単純tiltを有限回繰り返すと指定された反標準heartに到達する
（命題\ref{prop:numericaltorsion}，定理\ref{thm:Eulerchamber}）．
根類だけの比較でなく，実際のheartの一致を証明する点が重要である．

第(1)項はRainsの既知のt構造の記述の適用である．第(4)項は一般のcorner理論を
構造対象つき比較へ適用した判定法である．第(2)，(3)項は本研究の楕円gluingと共鳴の
結果をまとめたもので，必要な証明を本稿に収める．
\emph{全四heartの存在・一致・相対降下を一括して主張する定理ではない．}

さらに実際のrestriction条件\eqref{eq:two-nontrivial}を満たすAOU領域では，
最後の二つのmutationが対象を変えないことから
$E_0\simeq\pi F$を証明し，AS側の構造対象と反標準環を同定する
（定理\ref{thm:AOU-pointed}，系\ref{cor:AOU-ring}）．
これにより，corner有限性を確認した範囲ではAOUの指定関手自体のt-exactnessも得る
（系\ref{cor:AOU-texact}）．

\subsection{先行研究と本稿の位置}
Van den Bergh \cite{VdB}は非可換blow-upを構成し，section algebraへの比較を
一般位置条件の下で導来同値とした．Rains \cite{Rains,Filtered}はそのblow-upの
導来圏とt構造をgluingで記述し，effective rootと数値的なnef/ample coneを解析した．
その「ample」はeffective coneによる数値的概念であり，指定pairのAZ amplenessと
定義上同じではない．Rainsの2025年のmonograph \cite{RainsBook}は関連する先行稿を統合するが，
本稿の定理番号は実際に照合した2019年稿・2022年稿に従う．monographとの全番号対応は未確認である．

Abdelgadir--Okawa--Ueda \cite{AOU}は三次曲面型AS代数を構成し，
十分一般な楕円六点データについてVdBのsection modelとの導来比較を与えた．
Okawa--Ueda \cite{OU}はこのAS構成をdel Pezzo型へ拡張する．
これらを，raw blow-up heartとの無条件なabel同値とは読まない．
HRS tiltのrealizationが一般には同値でないことは\cite{CHZ}で詳しく論じられている．
本稿では既知の導来同値からtiltを記述する場合と，抽象的にtiltを作る場合を区別する．

VdBの\S11.2は既に境界非交叉対象の有限長性と単純対象数を扱う．
またADE型のheartとWeyl反射はKleinian特異点の安定性条件で既知であり
\cite{BridgelandKleinian}，単純tiltがlength heartを保つ一般条件は\cite{BCPW}で研究されている．
本稿はそれらの機構の発見を主張せず，反標準Euler関数による\emph{標的heartの選択}と，
VdB・Rains・AOUの\emph{指定関手}への適用を区別して述べる．

本稿の目的は，これらの構成を指定関手・構造対象・仮定つきで一つの比較体系にすることである．
既知のblow-up定理，HRS対応，AS定義自体を新しい結果とはしない．
AOUの構造対象の同定は，同論文のmutation列，境界上のadmissibility，
例外対のrestriction比較から導く明示的な帰結である．これをAOUの導来同値定理そのものの
新証明とはしない．本研究の組合せによる主張の網羅的新規性調査と外部独立査読は未了である．
また，境界作用を復元する数値的整合性検査を第\ref{sec:boundarydata}節に置く．
AOU v1の複数の表示を同じ境界同定で読むと非torsionの場合に不整合が生じるため，
その具体座標による辞書は未確定とし，pointed性を実際のrestriction条件の下で述べる．

\section{記法，realization，反標準環}\label{sec:setup}
基礎体は$k=\mathbb C$とし，圏は$k$-線形とする．導来比較には適合するdg enhancementを用いる．
$\cD$は冪等完備で，使用するRHomの全cohomologyは有限次元とする．
$\cH\subset\cD$が有界heartであるだけでは，realization
$\operatorname{real}_{\cH}:\Db(\cH)\to\cD$が同値とは限らない．
同値となる場合を本稿では\emph{derived type}と呼ぶ．

Serre関手を$\Sfun$，$\theta=\Sfun[-2]$，$s=\theta^{-1}=\Sfun^{-1}[2]$と書く．
$s$を単にAS代数の次数1 shiftと同一視しない．一般の$\mathbb Z$-代数には
一周期的なshiftの指定自体がなく，helixの一周期は複数頂点を通ることがある．
$s$がheartを両方向で保つなら
\[
\Ext^i_{\cH}(M,N)=\Hom_{\cD}(N,s^{-1}M[2-i])^*=0\quad(i>2)
\]
なので，derived typeの$\cH$のhomological dimensionは高々2である．
$\OO$がexceptionalなら$\Ext^2(\OO,s^{-1}\OO)=k$であり，ちょうど2となる．

\begin{definition}[構造対象つき比較]
$(\cH,\OO,s)$の間の\emph{pointed同値}とは，構造対象の同型と$s$との自然な両立同型を
指定した同値をいう．境界restrictionやmarkingの保存は，必要な箇所で別に指定する．
\end{definition}
構造対象を$e=\OO$と書き，右加群の規約で
\[
R=R(e,s)=\bigoplus_{n\ge0}\Hom(e,s^ne),
\qquad ab=s^{\deg b}(a)\circ b
\]
と置く．局所有限右noetherian性を確認した範囲で
$\AAC=\qgr R=\gr R/(\text{有限次元加群})$と定める．
$\AR$と$\AAC$を同じ三角圏のheartと呼ぶには，まず$\Phi$を構成してtransportする必要がある．

\begin{definition}[AZ ample pair]
noetherian abel圏のpair $(e,s)$がAZ-ampleであるとは，任意の対象が
有限個の$s^{-n}e$（$n>0$）の直和の商であり，任意の全射$M\twoheadrightarrow N$に対して
$\Hom(e,s^nM)\to\Hom(e,s^nN)$が$n\gg0$で全射となることをいう\cite{AZ}．
\end{definition}
高い全Extの消滅をこの二条件だけから無条件に推論しない．
またhelixの全頂点による生成性と，一つの頂点の$s$-orbitによる生成性は別である．

\section{比較の一般論：振幅1，ample orbit，corner}\label{sec:formal}
\subsection{振幅1の同値とHRS tilt}
\begin{proposition}\label{prop:hrs}
$F:\Db(\cA)\simeq\Db(\cB)$をenhancementと両立する同値，$G=F^{-1}$とし，
$F(\cA)\subset D^{[0,1]}(\cB)$とする．このとき
\[
\cT=\ker H^1F,\qquad\cF=\ker H^0F
\]
は$\cA$上のtorsion pairであり，
\[
G(\cB)=\{X:H^{-1}_{\cA}X\in\cF,\ H^0_{\cA}X\in\cT,
\ H^i_{\cA}X=0\ (i\ne-1,0)\}=\langle\cF[1],\cT\rangle.
\]
これはderived typeのheartである．$F$がt-exactであることと$\cF=0$は同値である．
\end{proposition}
\begin{proof}
有界t構造のaisleを有限拡大とshiftで記述すると，振幅仮定と直交性から
$G(\cB)\subset D^{[-1,0]}(\cA)$を得る．
$M\in\cA$，$B_i=H^i(FM)$に対しtruncation triangleへ$G$を施すと
$GB_0\to M\to GB_1[-1]$となる．長完全列から
$H^{-1}_{\cA}(GB_0)=H^0_{\cA}(GB_1)=0$．よって
\[
0\longrightarrow GB_0\longrightarrow M\longrightarrow
H^{-1}_{\cA}(GB_1)\longrightarrow0
\]
は$\cT$と$\cF$への分解である．$T\in\cT,U\in\cF$では
$FT\in\cB$，$FU\in\cB[-1]$なので$\Hom(T,U)=0$．
逆に二つのcohomologyが表示条件を満たす$X$は$FX\in\cB$となり，heartの式を得る．
このheartのrealizationは$G$を通して標準の$\Db(\cB)$と同定されるため同値である．
$\cF=0$ならtorsion分解から$\cT=\cA$で，両heartは一致する．逆は明らか．
\end{proof}
\begin{remark}
$\cF\ne0$の対象は，$F$で消えるのでなく$H^1(FM)[-1]\ne0$になる．
抽象的なtorsion pairからHRS heartを作るだけではderived type性は従わない\cite{CHZ}．
\end{remark}

\subsection{反標準orbitによるheartの一意性}
\begin{proposition}\label{prop:unique}
$\cH\subset\cD$はderived typeの有界heart，$s$は$\cH$を両方向で保ち，
$(e,s)$はAZ-ampleとする．さらに各$M\in\cH$について，すべての$p>0$に共通の閾値で
$\Ext^p(e,s^nM)=0$ $(n\gg0)$を仮定する．このとき
\[
\cH=\{X\in\cD:\Hom(e,s^nX[i])=0\quad(i\ne0,\ n\gg0)\}.
\]
従ってこの条件を満たすheartは固定した$(\cD,e,s)$に対し高々一つである．
\end{proposition}
\begin{proof}
有限個の$H^j_{\cH}X$に高次消滅を同時に適用し，hyper-Ext spectral sequenceを退化させると
$\Hom(e,s^nX[i])=\Hom(e,s^nH^i_{\cH}X)$となる．
非零$M$について右辺がすべての十分大きい$n$で零となることは，任意に大きい$N$で
$s^NM$を負twistの有限直和から生成するAZ条件に反する．従って指定した消滅は
$H^i_{\cH}X=0$ $(i\ne0)$と同値である．
\end{proof}
\begin{corollary}\label{cor:ampleexact}
二つのheartが命題\ref{prop:unique}の条件を満たすなら，pointed導来同値はt-exactである．
同じ条件のheartの非零$s$-固定対象$M$は$\chi(e,M)>0$を満たす．
$s^rM\simeq M$なら少なくとも$\sum_{a=0}^{r-1}\chi(e,s^aM)>0$である．
\end{corollary}
\begin{proof}
第一の主張はheartの認識式による．固定対象では高次Extが消え，
$\chi(e,M)=\dim\Hom(e,M)$となる．この値が零なら全orbitからのHomが零となってAZ生成性に反する．
周期対象には$\bigoplus_{a=0}^{r-1}s^aM$を用いる．
\end{proof}

\subsection{idempotent cornerの判定}
\begin{proposition}\label{prop:corner}
$B$は局所有限右noetherian非負次数付き代数，$p\in B_0$はidempotentとする．
$R=pBp$が右noetherian，$Bp$が右$R$有限生成，$pB$が左$R$有限生成であると仮定する．
自然な$(-)p:\qgr B\to\qgr R$が同値である必要十分条件は$\dim_k B/BpB<\infty$である．
\end{proposition}
\begin{proof}
$Bp$の有限生成性から$Mp$は有限生成であり，$(-)p$はexactかつ有限次元加群を保つ．
左随伴$L(N)=N\otimes_RpB$は有限生成加群を有限生成加群に送り，
$pB$の左有限生成性から有限次元加群も有限次元に送る．
$B/BpB$が有限次元なら，counit $Mp\otimes_RpB\to M$の核・余核は$p$で消される．
それらは右noetherian性により有限生成の$B/BpB$加群なので有限次元である．
$L$を短完全列へ施した際の左端の核も同じ議論で有限次元となり，$L$は商圏上exactとなる．
$(N\otimes_RpB)p=N$とcounitが逆同値を与える．
逆に$(-)p$が同値なら$\pi(B/BpB)$は零に送られるので零であり，有限生成性から有限次元である．
\end{proof}


\section{VdB heartとRains heart：blow-up表示の同定}\label{sec:VdBRains}
\subsection{blow-downからのt構造の認識}
Rains \cite[Lemma 4.9, Proposition 4.17]{Rains}はVdB blow-up
$\alpha:\widetilde X\to X$について，$X$とblow-up中心を含む可換divisorがGorensteinなら
$R\alpha_*$が\emph{pseudo-canonical}であることを示す．本稿で必要な内容は
\begin{equation}\label{eq:pseudocanonical}
M\in D^{\ge0}(\widetilde X)
\quad\Longleftrightarrow\quad
R\alpha_*\theta_{\widetilde X}^{\ell}M\in D^{\ge0}(X)
\quad(\ell\in\mathbb Z)
\end{equation}
という非負部分の認識である．初期heartとblow-down関手およびSerre作用を固定すれば，
右辺はt構造を一意に決める．

\begin{theorem}[同じblow-upデータに対する一致]\label{thm:RV}
VdBとRainsの二つの表示に共通の初期heartを取り，有限回の各blow-upについて，
中心，gluing bimodule，blow-down関手を同定する．各段階で
\eqref{eq:pseudocanonical}の適用条件が成り立つとする．
これらの同定から得るdg同値はt-exactで，$\AV\simeq\AR$を与える．
構造対象も同じラベルの対象として同定している場合，この同値はpointedである．
\end{theorem}
\begin{proof}
一段階について，基底のt-exact同値を$E$，上のdg同値を$F$とし，
$R\alpha_{\mathrm R*}F\simeq E R\alpha_{\mathrm V*}$を用いる．
Serre関手の一意性から$F\theta_{\mathrm V}\simeq\theta_{\mathrm R}F$である．
従って\eqref{eq:pseudocanonical}の右辺は$F$で保存され，非負部分が同定される．
直交性から非正部分も同定されるので$F$はt-exact．初期heartから帰納する．
構造対象の保存はt-exactnessからでなく，指定したgluingラベルの同定から従う．
\end{proof}
\begin{remark}
Rains \cite[Theorem 4.10, Proposition 4.13, Theorem 4.30]{Rains}は
VdB blow-upのsemiorthogonal decompositionとgluingを用いる．従って同じ構成を
この二通りに記述した場合には上の仮定を満たす．しかし，独立に選んだ
blowdown構造の比較，異なるパラメータ表示の照合，途中でsection modelに置き換える操作は，
同じデータという言葉だけで済ませられない．
\end{remark}

\subsection{section modelへの移行は別の操作}
VdB \cite[\S11.1]{VdB}の図式を
\[
X_j\xleftarrow{\alpha_j}\widetilde X_j\xrightarrow{\delta_j}X_{j+1}
\]
とする．$\alpha_j$はblow-up，$\delta_j$はsection-algebra modelへの比較である．
Definition 11.1.2の一般位置条件を，同論文で固定した平面埋込みに関する
$(\sigma p_i)$へ課すと，Theorem 11.1.3より$R\delta_{j*}$と$L\delta_j^*$は逆同値で，
heart上の振幅は$[0,1]$と$[-1,0]$．coherent有界部分への制限にはLemma 9.1.5も用いる．
命題\ref{prop:hrs}から
\[
\cT_j=\ker R^1\delta_{j*},\qquad \cF_j=\ker\delta_{j*},\qquad
(L\delta_j^*)(\coh X_{j+1})=\langle\cF_j[1],\cT_j\rangle.
\]
$\cF_j=0$のとき，かつそのときに限り，この比較はabel同値である．

三次曲面では$X_7=\operatorname{Proj}F$である．VdBのChapter 12の表示$F=P/(C_3)$は，
Theorem 11.1.3の一般位置条件を仮定せずにも構成される．
従って超曲面表示の存在からその点での導来同値を結論してはいけない．
同書Remark 12.0.5も，raw blow-upとsection modelが非可換の場合に一般には一致しないことを区別する．
六段の比較の合成を，対象・関手の照合なしに一回のtiltと呼ぶこともしない．

\section{Rains heartと反標準heartの導来比較}\label{sec:Rac}
\subsection{genuine楕円gluingと既知入力}
$S$を滑らかな可換del Pezzo曲面で$-K_S$がampleなものとし，
$E\in|-K_S|$を滑らかな楕円曲線，$q\in\Pic^0(E)$とする．
$\cB=\OO_S^\perp\subset\Perf(S)$にexceptional object $e$を
\begin{equation}\label{eq:glue}
\RHom(e,M)=0,\qquad \RHom(M,e)=\RHom_E(Li^*M,q)
\quad(M\in\cB)
\end{equation}
でglueしたdg圏を$\cD_q$とする．Rainsの構成から得るheartを$\AR$とし，
$\cD_q=\Db(\AR)$とする．これは任意の抽象的gluingにheartの存在を仮定するものではなく，
\cite[Theorem 4.30, \S5]{Rains}と\cite[\S4]{Filtered}の構成を指定したものである．
その境界関手$F_q$と右随伴$i_{q*}$は
\[
F_q(e)=\OO_E,\quad F_q|_{\cB}=Li^*(-)\otimes q^{-1},\quad
\theta X\longrightarrow X\longrightarrow i_{q*}F_qX\longrightarrow\theta X[1]
\]
を満たす．$s=\theta^{-1}$はheartを保ち，$F_qs=\Psi F_q$で，
$\Psi$は境界の平行移動と次数$d=(-K_S)^2$の線束によるtensor積の合成である．

Rains \cite[Proposition 5.1]{Filtered}を$e$へ適用すると
\begin{equation}\label{eq:ring}
R=\bigoplus_{n\ge0}\Hom(e,s^ne),\quad g\in R_1\text{ は中心正則},\quad
R/gR=B(E,L,\sigma),\quad\deg L=d>0.
\end{equation}
ここで$B(E,L,\sigma)$はtwisted homogeneous coordinate ringである．
$R$は$A=R/(g-1)$のRees環であり，$\operatorname{gr}_{(g)}R=B(E,L,\sigma)[g]$．
固定次数での$g$進filtrationは有限なので，整域性と左右noetherian性は商から持ち上がる．
$\sigma$が無限位数なら\cite[Proposition 2.4]{RSS}により$\chi$条件と有限cohomological dimensionを得る．
この命題は$\deg L=1$も許す．有限位数なら$B(E,L,\sigma)$はskew Laurent環の部分環としてPIであり，
\cite[Theorems 5.13, 5.1(1), 4.5]{YZ}を順に用いると同じ有限性を得る．
ここで必要なのは$\pi R(a)$からのExtの有界性で，標的の有限大域次元を先取りしてはいない．

\begin{lemma}[orbit間のExt]\label{lem:orbit}
すべての整数$a,b$に対し，sourceの$s^ae,s^be$とtargetの$\pi R(a),\pi R(b)$の双方で
\begin{equation}\label{eq:orbit}
\Ext^j(-,-)=
\begin{cases}
R_{b-a}&j=0,\ b\ge a,\\
R_{a-b-1}^*&j=2,\ b<a,\\
0&\text{その他}
\end{cases}
\end{equation}
となる．targetの次数2のpairingは境界$H^1(E,\OO_E)$へのtraceで定まる．
\end{lemma}
\begin{proof}
sourceでは$e$のexceptionalityと境界triangleから始め，
正次数の楕円線束の$H^1$消滅で帰納する．$g$の単射性とSerre双対性が負のtwistを処理する．
targetの消滅とtrace pairingは\cite[Corollary 3.8とその証明]{Filtered}である．
Homの同定が自然な$R_n\to\Hom(\pi R,\pi R(n))$であることは，
$g$の短完全列に対する切断の列と$n=0$での$\End(\pi R)=k$から帰納して確認する．
\end{proof}

\subsection{orbit生成，trace，飽和性}
\begin{theorem}[指定した自然導来切断の同値]\label{thm:derived}
上のgenuine楕円gluingに対し
\[
\Phi=\pi\mathbf R\Gamma_*:\cD_q\xrightarrow{\sim}\Db(\qgr R),
\qquad (\mathbf R\Gamma_*X)_n=\RHom(e,s^nX)
\]
はdg同値である．$\Phi(e)=\pi R$，$\Phi s\simeq(1)\Phi$を満たし，
$\qgr R$のhomological dimensionは2である．$q$と$\sigma$の位数に制限はない．
\end{theorem}
\begin{proof}
\textit{第一段階：Ind圏での生成．}
$X\in\Ind(\cD_q)$がすべての$s^ne$に右直交するとする．
$\theta X$も同じ条件を満たすので，境界triangleと随伴性から
\[
\RHom_E(\Psi^n\OO_E,F_qX)=0\quad(n\in\mathbb Z)
\]
となる．正次数のtwisted line bundlesは$D_{\mathrm{qc}}(E)$をcompact生成するから$F_qX=0$．
さらに$\RHom(e,X)=0$より$X\in\Ind(\cB)$．対応する$G\in D_{\mathrm{qc}}(S)$は
$R\Gamma(S,G)=0$，$Li^*G=0$を満たす．$U=S\setminus E$はaffineである．
$j:U\hookrightarrow S$に対しCartier divisorの局所化から
\[
G\simeq\operatorname*{hocolim}_{n}G(nE)\simeq Rj_*j^*G.
\]
affine上の大域切断の保守性より$G=0$，従って$X=0$．これでorbitはInd圏でcompact生成し，
compact objectsは$\cD_q=\thick\{s^ne:n\in\mathbb Z\}$となる．

\textit{第二段階：関手と完全忠実性．}
Rains heartのInd圏のinjective resolutionを取り，exactな$s$の作用をstrictifyして
右$R$作用つき$\mathbf R\Gamma_*$を構成する．まず$D(\QGr R)$への関手である．
\eqref{eq:orbit}より$\Phi(s^ae)=\pi R(a)$で，高次cohomologyは有限内部次数にのみ残って$\pi$で消える．
thick生成により全対象の像は$\Db(\qgr R)$へ入る．
境界の$V\in\coh E$については$R\Gamma(E,\Psi^nV)$となり，Serre消滅により
$\Phi(i_{q*}V)$は通常の境界包含$\coh E\simeq\qgr(R/gR)\to\qgr R$の像である．
この同定は短完全列を保存する．

非零の$u\in H^1(E,\OO_E)$を$\OO_E$の自己拡大で表す．境界へ押し出し，
$e\to i_{q*}\OO_E$で引き戻し，境界triangleのconnecting mapを取ると
$\Ext^2(e,\theta e)$の非零元を得る．$e$のexceptionalityから
$H^1(E,\OO_E)\simeq\Ext^2(e,\theta e)$である．targetでも同じ議論が成り立ち，
$\Phi$は$u$を同じ$u$へ送る．従ってtraceの写像は非零．
全shiftとYoneda積に対する自然性から，orbit間の次数2の写像は次数0の同型の双対となる．
残る次数は\eqref{eq:orbit}で消えるので，orbit間のRHomはquasi-isomorphism．
両変数でthick closureを取れば$\Phi$の完全忠実性が従う．

\textit{第三段階：本質的全射性．}
$\cB$は$\Perf(S)$のadmissible部分圏でsmooth properである．gluing moduleは
pseudo-perfectであり，smoothnessからperfectとなる．従って\cite[Theorems 3.18, 3.25]{Orlov}より
$\cD_q$はsmooth properかつ飽和的である．任意の$Y\in\Db(\qgr R)$に対し
$X\mapsto\RHom(\Phi X,Y)$は，環の有限性とorbit生成からpseudo-perfect moduleとなり，
sourceの対象で表現される．こうして右随伴$\Phi^!$を得る．
counitのcone $K$はすべての$\pi R(a)$に右直交する．
標的の標準ample pair，$\chi$と有限cohomological dimensionを使い，有限個の$H^iK$へ
同時にSerre消滅を適用すると，非零cohomologyがあれば高いtwistの非零大域切断を得て矛盾する．
よって$K=0$であり$\Phi$は本質的全射．

同値でSerre関手を運ぶとtarget上で$(-1)[2]$となる．$(-1)$はheartを保つので
Extは次数2を超えて消える．$\Ext^2(\pi R,\pi R(-1))=k$からhomological dimensionはちょうど2である．
\end{proof}
\begin{remark}
この証明は共鳴により元のheartのAZ amplenessが失敗しても働く．
ただし$S\setminus E$がaffineであることを本質的に使う．
weak del Pezzo，特異楕円境界，非局所自由gluingにそのまま適用はしない．
また全$q$で同値という表現は，各$q$のsourceと対応するtargetの比較を意味する．
\end{remark}

\subsection{全heartの一回のtilt}
\begin{theorem}\label{thm:globaltilt}
定理\ref{thm:derived}の設定で$H=-K_X$がnefなRains heartを取る．
全$M\in\AR$に対し$\Phi(M)\in D^{[0,1]}(\AAC)$であり，
\[
\cT=\ker H^1\Phi,\quad\cF=\ker H^0\Phi,\quad
\Phi^{-1}(\AAC)=\langle\cF[1],\cT\rangle.
\]
\end{theorem}
\begin{proof}
Rains \cite[Theorem 9.10]{Rains}によるline bundlesの生成性を用い，
$\bigoplus_j L_j\twoheadrightarrow M$を取る．Serre双対性から
\[
\Ext^2(e,s^nM)=\Hom(M,\theta^{n+1}e)^*.
\]
このHomは$\bigoplus_j\Hom(L_j,\theta^{n+1}e)$へ単射する．
非零射があれば\cite[Corollary 10.12]{Rains}により
$(n+1)K_X-c_1(L_j)$が有効だが，その$H$との交点は
$-(n+1)d-Hc_1(L_j)<0$ $(n\gg0)$でありnef性に反する．
従って次数2のcohomologyは有限内部次数にしかなく，$\pi$で消える．
負次数と次数$>2$はheart性とhomological dimensionで消える．命題\ref{prop:hrs}を適用する．
\end{proof}
\begin{remark}
$\cT,\cF$は$H^j\Phi$によって定義される．境界外でない一般の対象について，
これを生の$\Ext^j(e,M)$の消滅と同一視してはならない．
$\Phi$のt-exactnessは正確に$\cF=0$と同値である．
t-exactならpointed性により$(e,s)$はAZ-ampleとなる．逆向きの認識には，
例えば命題\ref{prop:unique}の高次消滅条件も確認すれば十分である．
\end{remark}

\section{共鳴，数値的ampleness，AZ ampleness}\label{sec:resonance}
\subsection{共鳴の向きと実在条件}
古典的avatarの根格子$K_S^\perp$の根を$\alpha^2=-2$，$\alpha K_S=0$で定める．
$\rho(\alpha)=\OO_E(\alpha)$と書き，rank-zeroの類について
\begin{equation}\label{eq:resonance}
Li^*M=0\quad\Longrightarrow\quad \rho(\alpha)\otimes q^m\simeq\OO_E,
\qquad\alpha=c_1(M),\quad m=\chi(e,M)
\end{equation}
となる向きに$q$を取る．$q$が非torsionなら$m$は一意である．
これはrestrictionのdeterminantの条件であり，特定のheartのsheafの存在とは別である．
特に$\alpha$と$-\alpha$を入れ替えるだけで，実在sheafの指数を自由に反転することはできない．

Rainsの数値的ample coneは有効divisor錐の双対の内部である\cite[\S4]{Filtered}．
有効根$\alpha$があれば$H\alpha=0$なので$H$はこの意味ではampleでない．
以下では，そのことと指定pair $(e,s)$のAZ amplenessの失敗を別に証明する．

\begin{theorem}[既約根sheafの符号別Ext]\label{thm:sign}
境界非交叉の既約純一次元root sheaf $M$が実在するとする．$m=\chi(e,M)$なら全$n\in\mathbb Z$で
\[
\dim\Hom(e,s^nM)=\max(m,0),\quad
\dim\Ext^1(e,s^nM)=\max(-m,0),\quad\Ext^2(e,s^nM)=0.
\]
$M$は2-sphericalであり，自然導来切断の像は
\[
\Phi(M)\simeq
\begin{cases}
V_m & m>0,\\ 0&m=0,\\V_{-m}[-1]&m<0,
\end{cases}
\]
ただし$V_{|m|}$は境界外の$|m|$次元$A=R/(g-1)$加群に対応する対象である．
\end{theorem}
\begin{proof}
境界triangleから$\theta M\to M$は同型なので$s^nM\simeq M$．
line bundleから一次元以下のsheafへのExt$^2$は\cite[Corollary 7.23]{Rains}で消える．
$m\le0$で非零$e\to M$があれば，像$I$は$c_1(I)=\alpha$を持ち，余核$T$は零次元．
$I$は$e$の商なので\cite[Corollary 10.16]{Rains}より
$\chi(e,I)\ge-\alpha(\alpha+K_X)/2=1$．零次元$T$のEuler値は非負なので$m\ge1$となり矛盾する．
従ってこの場合$h^0=0$，$h^1=-m$である．

$m>0$ではRainsの反変導来双対を使う．純一次元sheafの双対は
$\mathbf R\mathrm{ad}(M)=N[-1]$の形で，$N$も双対曲面の既約純一次元境界外根sheafである
\cite[Corollaries 5.9, 7.26, 7.30, 7.31]{Rains}．双対構造line bundleを$e^{\mathrm{ad}}$とすると
\[
\Ext^j(e^{\mathrm{ad}},N)\simeq\Ext^{j+1}(M,e)
\simeq\Ext^{1-j}(e,M)^*\quad(j=0,1).
\]
従って$\chi(e^{\mathrm{ad}},N)=-m$．前段の議論を双対側の指定line bundleへ適用して
$h^1(e,M)=0$，$h^0(e,M)=m$を得る．

既約性と純粋性から非零自己射は単射で，同じ数値類の間の単射は同型である
\cite[Corollary 7.7]{Rains}．従って$\End(M)=k$．Serre双対性と
$\chi(M,M)=-\alpha^2=2$から自己Extは次数0,2で$k$，その他で零．
最後に$g$は$M$上で可逆に作用するから，非零の次数付きExt加群は
$k[g]$上有限自由で，局所化の次数0部分が所要の有限次元$A$加群となる．
\end{proof}

\begin{corollary}\label{cor:negative}
$m<0$の上記$M$が存在すれば$(e,s)$はRains heartでAZ-ampleでなく，$\Phi$はt-exactでない．
$\Phi$でtransportした反標準heartはRains heartと一致せず，両者の標準pointed abel圏の同値もない．
定理\ref{thm:derived}の範囲では，それでも導来同値は成立する．
\end{corollary}
\begin{proof}
$\Hom(e,s^nM)=0$が全$n$で成り立つためAZ生成性に反する．
$\Phi(M)=V[-1]$は非零であり，標準heartに属さない．pointed abel同値はample pairを保存する．
最後の主張は定理\ref{thm:derived}による．
\end{proof}
\begin{remark}[三つの非同値な条件]
数値的共鳴，実在する根sheaf，AZ amplenessの失敗は同義語でない．
正指数の個々の根sheafは$\cT$に入り，それだけではAZ失敗を検出しない．
負指数の根sheafが見つからないことだけから$\cF=0$とも結論しない．
またheartの非一致は指定した同値の下での主張であり，構造対象を忘れた任意のabel同値の不存在ではない．
\end{remark}

\subsection{稠密性と零指数の境界}
\begin{proposition}\label{prop:density}
滑らかな平面三次$E$，非torsion $q$と，一般の五点$p_1,p_2,p_4,p_5,p_6\in E$を固定する．
$L=\OO_E(1)$とし，$p_3^{\varepsilon}(n)$を
\[
\OO_E(p_3^{\varepsilon}(n))=L(-p_1-p_2)\otimes q^{-\varepsilon n},
\qquad n\ge1,\quad\varepsilon=\pm1
\]
で定める．各符号で有限個を除き六点は一般位置であり，
$h-e_1-e_2-e_3$の指数は$m=-\varepsilon n$となる．
各符号の点列は$E$にZariski稠密である．既約sheaf存在結果が適用される範囲で，
負指数の列はAZ失敗の稠密な例を与える．
\end{proposition}
\begin{proof}
$\Pic^1(E)\simeq E$より点は一意で，$q$の無限位数性から相異なる．
固定五点に対し，三点共線または六点conicの条件を生む追加点は有限個である．
残る無限集合は曲線上稠密で，restrictionの計算が指数を与える．
\end{proof}
\begin{proposition}\label{prop:zeroroot}
$S$がgenuineで$E\in|-K_S|$が滑らかなら，非零根$\alpha$について
$\OO_E(\alpha)\not\simeq\OO_E$である．
\end{proposition}
\begin{proof}
$H=-K_S$のamplenessから$H^0(\pm\alpha)=0$．
$\chi(\alpha-H)=0$，$H^0(\alpha-H)=H^2(\alpha-H)=0$なので$H^1(\alpha-H)=0$．
$0\to\OO_S(\alpha-H)\to\OO_S(\alpha)\to\OO_E(\alpha)\to0$から
$H^0(E,\OO_E(\alpha))=0$となる．
\end{proof}
弱del Pezzoの$(-2)$曲線$D$では，対照的に$i_*\OO_D(-1)$が全反標準orbitに直交し，
反標準収縮で導来的に消える．これは$m=0$の現象であり，非零負指数のshiftと区別する．
実際，$H|_D=\OO_D$なので全$n$で$R\Gamma(D,\OO_D(-1)\otimes\OO_D(nH))=0$である．


\section{AS heartと反標準heart：helixからcornerへ}\label{sec:AS}
\subsection{AS構成が指定するもの}
AS heartは，適切なnoetherian性を満たす三次元Artin--Schelter正則
$\mathbb Z$-代数の$\qgr$として定義する．三次曲面の場合には
$\{0,1,2\}\times\mathbb Z$を添字とする代数を使い，標準的な九対象のfoundationがある
\cite[\S3--4]{AOU}．より一般のdel Pezzo型の標識とAS条件は
\cite[Definitions 1.1, 3.1]{OU}に従う．
ここでのAS正則次元3は曲面のhelix代数の次元であり，
非可換射影三空間を与える四次元AS正則ambientと同じものではない．
また非周期的な$\mathbb Z$-代数を，根拠なく一周期の次数付き代数に置き換えない．

very strong helixは例外性，特定の周期範囲でのstrong性，Serre作用との周期関係を与える．
そのhelixがheartを再構成するためには，生成性と高次Ext消滅を伴う豊富性も用いる．
しかし，九対象の全orbitの豊富性から，そのうちの一つ$e$のorbitの豊富性は自動的には従わない．
本稿で反標準環を作るのは後者であり，ここにcornerを導入する理由がある．

\subsection{既知の導来同値と構造対象の課題}
AOU \cite[\S9]{AOU}は，滑らかな楕円曲線，次数3の線束，無限位数の平行移動$\sigma$と，
sufficiently generalな六点を用いる．同節では最初から
$\coh\widetilde X=\qgr F$と置いている．この$\widetilde X$をVdBのraw blow-upと
定義上同一視してはならない．Lemma 9.2のSerre関手は$(-1)[2]$であり，
Theorem 9.4はAS代数$\mathsf A$に対するdg導来同値
\begin{equation}\label{eq:AOU}
\Psi:\Db(\qgr F)\xrightarrow{\sim}\Db(\qgr\mathsf A)
\end{equation}
を与える．両側のfull strong collectionのEnd代数を同定する証明であり，
これは標準heart間のabel同値を主張する定理ではない．

mutation後のcollectionの先頭を$E_0\in\Db(\qgr F)$とすると，
同値は$\Psi(E_0)\simeq\OO_{\mathrm{AS}}(0,0)$と規格化できる．
\cite[(9.44)]{AOU}の可換表示では先頭は$\OO$だが，その後に非可換版の対象を
mutationによって定義し直している．従ってpointed比較には
\begin{equation}\label{eq:E0}
E_0\stackrel{?}{\simeq}\pi F
\end{equation}
を実際の対象として確認する必要がある．$K_0$類や楕円境界へのrestrictionの一致では足りない．
この同定を未確認の範囲では，$R(E_0,s)$と$R(\pi F,(1))$を同じ環と呼ばない．
一方，反Serre作用との両立性自体はSerre関手の一意性から得られる．

\begin{corollary}[VdB--ASの導来比較の範囲]\label{cor:VdBAS}
VdBの最後のsection-model比較とAOUの比較\eqref{eq:AOU}の両方が適用でき，
中間の$\qgr F$を同じデータで同定しているなら，それらの合成は
VdBの当該blow-up heartとAS heartの導来同値を与える．
そのt-exactnessと構造対象の保存は，これだけでは従わない．
\end{corollary}
\begin{proof}
二つの指定dg同値を合成する．VdB側では既に非自明なHRS tiltがあり得るため，
合成がt-exactであることは同値性だけの帰結ではない．構造対象の保存には
各関手での像，特に\eqref{eq:E0}の確認が必要である．
\end{proof}

\subsection{一対象の反標準orbitを判定する}
AS側の標識されたfoundationを$E_1=e,E_2,\ldots,E_r$とし，
$T_0=\bigoplus_iE_i$，$B_n=\Hom(T_0,s^nT_0)$と置く．
$E_1$への射影を$p\in B_0$とすると，積の規約を合わせれば
\begin{equation}\label{eq:foldedcorner}
pBp=\bigoplus_{n\ge0}\Hom(e,s^ne)=R(e,s).
\end{equation}
以下の条件は，任意のAS退化点で自動的に成立するものではない．

\begin{corollary}[AS--反標準のabel同値判定]\label{cor:ASac}
上のfoldingが$\AAS\simeq\qgr B$を与え，$e$が$\pi(pB)$に対応するとする．
$B$と$R=pBp$が局所有限右noetherianで，$Bp$が右$R$有限生成，
$pB$が左$R$有限生成であると仮定する．このとき自然なcorner関手
\[
K:\AAS\simeq\qgr B\longrightarrow\qgr R(e,s)
\]
は，$\dim_k B/BpB<\infty$と同値にabel同値となる．
その場合$K(e)=\pi R$であり，確認したfoldingの下で$s$と$(1)$にも適合する．
\end{corollary}
\begin{proof}
式\eqref{eq:foldedcorner}と命題\ref{prop:corner}を使う．
$(pB)p=pBp$が構造対象の保存を与える．反標準次数のshiftと$s$の適合性は
foldingの定義に含めたものである．
\end{proof}

この系が同定するのは，\emph{AS側の指定対象から作った}反標準heartである．
これをRains側の$R(e,s)$やAOUの$F$と同定するには，pointedな比較を別に必要とする．
$B/BpB$が無限次元でも，それはこの自然なabelian corner関手が同値でないことを意味するだけで，
別の導来同値の不存在ではない．

\begin{lemma}[有限次数による判定が有効になる条件]\label{lem:cutoff}
$C=B/BpB$が非負次数付きで，有限次元$C_0$上，次数1で生成されるとする．
このとき$C_N=0$なら全$n\ge N$で$C_n=0$である．
従って$C$が有限次元であることと，ある$N$で$C_N=0$であることは同値である．
\end{lemma}
\begin{proof}
任意の次数$n\ge N$の生成単項式は，長さ$N$の先頭部分を含む．
$C_N=0$ならそれが零となる．逆は次数ごとの局所有限性による．
\end{proof}
三次曲面のJacobi表示では道長次数と反標準cut次数は異なる．
補題を道長次数に適用した結果を反標準cornerへ移すには，この二つのgradingの関係を確認する．
単に有限次数の計算が良好だったという理由だけで，全次数の有限性を結論しない．

\begin{remark}[三次曲面の具体的corner計算との関係]
同プロジェクトの別稿\nolinkurl{cubic_heart_corner.tex}は，指定された九頂点表示について，
四交比の十一のdeterminant条件，一様な道長7での消滅，および
$8+21t+33t^2+38t^3+33t^4+21t^5+8t^6$という有限cornerのHilbert関数を扱う．
それらの詳細な幾何学的証明と計算資産は別稿に保持する．本稿はその数値を
全AS表示への無条件な認識定理として用いず，系\ref{cor:ASac}の有限性条件の
具体的な検査法として位置づける．
\end{remark}


\section{混合\texorpdfstring{$A_2$}{A2}によるheartのずれの具体化}\label{sec:A2}
この節は定理\ref{thm:derived}のgenuine楕円gluingで，$q$非torsionとする．
$H=-K$はnef，$H^2>0$．楕円境界と交わらない一次元以下の対象を
$\cZ=\{M\in\AR:\dim M\le1,\ F_qM=0\}$と書く．

\subsection{単純根sheafと有限長部分圏}
\begin{lemma}\label{lem:offboundary}
$\cZ$は有限長Serre部分圏であり，その既約純一次元root sheafは$\AR$の単純対象である．
\end{lemma}
\begin{proof}
境界非交叉な対象の零次元商$Z$には$i^*Z=0$が従い，$\theta Z\to Z$は全射となる．
両者の数値類は等しいので\cite[Corollary 7.7]{Rains}から同型である．
$[Z]=\ell[\mathrm{pt}]$，$\ell\ge0$に対して境界determinantは$q^{\ell}$であり，
非torsion性から$\ell=0$，従って$Z=0$．既約純一次元sheafの非零部分対象の商は
零次元なので，この場合には単純性が得られる．

$0\to N\to M\to Q\to0$，$M\in\cZ$とすると，nef性と加法性から
$Hc_1(N)=Hc_1(Q)=0$．$i^*Q=0$なので$\theta Q\to Q$は全射であり，
rank零，$Kc_1(Q)=0$と\cite[Corollary 7.15]{Rains}から数値類は変わらない．
Corollary 7.7と自然性により$\theta Q\to Q$，$\theta N\to N$はいずれも同型で，
$N,Q\in\cZ$となる．拡大閉性は境界triangleから従う．
真の降鎖の逐次商は零次元になれないため，\cite[Proposition 7.32]{Rains}の
finitely partitive性が降鎖を止める．coherent heartのnoetherian性と合わせて有限長となる．
\end{proof}

\begin{corollary}[全境界外圏のtorsion pair]\label{cor:Ztilt}
定理\ref{thm:globaltilt}のtorsion pairは$\cZ$へ制限され，
\[
\cT\cap\cZ=\ker\Ext^1(e,-),\qquad
\cF\cap\cZ=\ker\Hom(e,-)
\]
となる．これは二根部分圏だけでなく全$\cZ$についての結論である．
\end{corollary}
\begin{proof}
$\cZ$がSerre部分圏なので，その対象のtorsion分解の部分対象と商も$\cZ$に属する．
$M\in\cZ$では$g$が可逆に作用し，各$\Ext^j(e,s^nM)$の移行写像が同型となる．
その局所化を$V^j(M)$とすれば$\dim V^j(M)=\dim\Ext^j(e,M)$であり，
$H^j\Phi(M)=0$と$\Ext^j(e,M)=0$は同値である．
\end{proof}

異なる二根sheaf$S_\alpha,S_\beta\in\cZ$には
\begin{equation}\label{eq:rootext}
\Hom(S_\alpha,S_\beta)=\Ext^2(S_\alpha,S_\beta)=0,
\qquad \dim\Ext^1(S_\alpha,S_\beta)=\alpha\cdot\beta
\end{equation}
が成り立つ．実際，単純性からHomが消え，境界非交叉とSerre双対性からExt$^2$も消える．
rank-zero Euler形式$\chi(S_\alpha,S_\beta)=-\alpha\beta$が残る次元を与える．

隣接根$\alpha\beta=1$，$\chi(e,S_+)=p>0$，$\chi(e,S_-)=-n<0$を固定し，
両向きの非分裂拡大を
\begin{equation}\label{eq:PQ}
0\to S_-\to P\to S_+\to0,
\qquad 0\to S_+\to Q\to S_-\to0
\end{equation}
とする．存在は\eqref{eq:rootext}から従い，各中間項は同型を除き一意である．

\begin{proposition}[connecting mapの最大rank]\label{prop:rank}
$P$のconnecting map$\delta:\Hom(e,S_+)\to\Ext^1(e,S_-)$は
$\rk\delta=\min(p,n)$を満たす．従って
\[
(h^0(P),h^1(P))=((p-n)_+,(n-p)_+),\qquad (h^0(Q),h^1(Q))=(p,n),
\]
ただし$h^j(M)=\dim\Ext^j(e,M)$，$x_+=\max(x,0)$とする．
このrankの証明には定理\ref{thm:derived}を必要としない．
\end{proposition}
\begin{proof}
符号別Extから$\Hom(e,P)=\ker\delta$，$\Ext^1(e,P)=\operatorname{coker}\delta$．
非零$f:e\to P$があれば，$P\to S_+$との合成は非零で全射となる．
像と$S_-$との交わりは零か$S_-$であり，前者なら拡大が分裂するので後者である．
従って$f$は全射．核$L$はrank-one torsion-freeで$c_1(L)=-(\alpha+\beta)$．
Corollary 10.12と$H(K+\alpha+\beta)=-H^2<0$により
$\Ext^2(e,L)=\Hom(L,\theta e)^*=0$．$e$のexceptionalityから$h^1(P)=0$となる．
$p>n$なら$\ker\delta\ne0$なので$\delta$は全射．$p\le n$で$h^0(P)>0$とすると
$\chi(e,P)=h^0(P)>0$に矛盾するので，この場合は単射である．
$Q$には逆向き短完全列の六項完全列を使う．
\end{proof}

\begin{proposition}[二根の全有限長対象]\label{prop:A2category}
$\cC=\langle S_+,S_-\rangle_{\mathrm{Serre}}$は
$\operatorname{mod}_{\mathrm{fd}}\Pi(A_2)$と同値であり，直既約対象は$S_+,S_-,P,Q$の四つである．
ここで$\Pi(A_2)$は二頂点・両向き二矢と，両方の長さ二の道を零とする関係を持つ代数である．
$p>n$では$(\cT\cap\cC,\cF\cap\cC)=(\add(S_+,P),\add(S_-))$，
$p<n$では$(\add(S_+),\add(S_-,P))$である．
\end{proposition}
\begin{proof}
非分裂性と単純性から，順序$(S_+,S_-,P,Q)$でHom次元は
\[
\begin{pmatrix}
1&0&0&1\\0&1&1&0\\1&0&1&1\\0&1&1&1
\end{pmatrix}.
\]
例えば$\Hom(P,Q)$は$P\twoheadrightarrow S_+\hookrightarrow Q$で生成され，
逆方向も同様で，両合成は零である．$\theta$はこれらの対象で自然に恒等なので
\[
\dim\Ext^1(U,V)=\dim\Hom(U,V)+\dim\Hom(V,U)+c_1(U)c_1(V).
\]
上の表を代入すると，$P,Q$から二単純へのExt$^1$はすべて零である．
長さによる帰納法から$P\oplus Q$は$\cC$のprojective generatorであり，
そのEnd代数は表示した四次元代数となる．
その表現$x:U_+\to U_-$，$y:U_-\to U_+$，$xy=yx=0$は，
$\operatorname{im}y\subset\ker x$，$\operatorname{im}x\subset\ker y$に適合する基底により，
一方の矢だけが非零の一次元対と孤立一次元空間へ分解される．従って直既約は四つである．
最後の主張は命題\ref{prop:rank}と符号別Extによる．
\end{proof}

$\cC$の内在的ExtはambientのExtと異なる．実際$\Pi(A_2)$の単純加群には
二周期的なprojective resolutionがあり，$\Ext^{2j}_{\cC}(S_+,S_+)=k$ $(j\ge0)$である．
従って$\Db(\cC)\to\Db(\AR)$は完全忠実でない．
また$p=n$なら$P$は自然な導来切断で消えるが，genuine avatarでは
$\alpha+\beta$が零指数根となるためこの同時共鳴自体が排除される．

\subsection{反標準側の二つのcharacter}
\begin{example}[実在する指数$(2,-1)$]
滑らかな平面三次$E$の変曲点を原点に取り，非torsion $q$について
\[
(p_1,\ldots,p_6)=(q,3q,2q,5q,7q,11q)
\]
とする．三点共線なら対応する正整数の和の$q$倍が零となり，六点conicも同様なので，
六点は一般位置である．$\alpha=e_1-e_2$，$\beta=e_2-e_3$の指数は$2,-1$．
Rainsのsimple-root effectivityと既約sheaf存在結果
\cite[\S4]{Filtered}\cite[\S9,\S11]{Rains}がこのchartで適用されると，
上記の$S_+,S_-$が同じheartに実在する．
数値計算だけをsheaf存在の代わりには用いない．
\end{example}

\begin{lemma}[境界外でのExtの比較]\label{lem:localization}
有限次元$A=R/(g-1)$加群$U,V$を境界外の$\qgr R$の対象とみなすと，
全$j\ge0$で$\Ext_A^j(U,V)\simeq\Ext_{\qgr R}^j(U,V)$である．
\end{lemma}
\begin{proof}
$R[g^{-1}]$は強次数付き環で，その次数付き加群圏は$\operatorname{Mod}A$と同値である．
exactな中心局所化の右随伴はrestrictionで，これもexactかつ完全忠実である．
局所化加群は内部次数torsionに導来的に直交するため，右随伴は$\QGr R$にも降りる．
よって導来随伴も完全忠実となる．有限次元加群の非負次数格子は$k[g]$上有限自由で，
有限生成$R$加群を与える．locally noetherian圏でのcoherent Extとの一致により主張を得る．
\end{proof}

\begin{theorem}[三角表現と四次元商]\label{thm:affinequotient}
上の指数$(2,-1)$の例で$V=\Phi(S_+)$，$U=\Phi(S_-)[1]$，$W=\Phi(P)$とする．
これらはそれぞれ次元$2,1,1$の$A$加群であり，
\[
0\longrightarrow W\longrightarrow V\longrightarrow U\longrightarrow0
\]
は非分裂で，$W\not\simeq U$である．$A$の$V$への作用像は全三角行列代数となる．
逆向き非分裂拡大$0\to U\to V'\to W\to0$も存在し，
\[
A/\Ann_A(V\oplus V')\simeq\Pi(A_2).
\]
新しい二単純の根類は$\alpha+\beta,-\beta$であり，指数は$1,1$である．
\end{theorem}
\begin{proof}
最大rankから短完全列を得る．定理\ref{thm:derived}，補題\ref{lem:localization}と
$\Ext_X^1(P,S_-)=0$から$\Hom_A(W,U)=0$となり，二つの一次元加群は異なる．
非零包含$S_-\to P$は同値によって非零の$\Ext_A^1(U,W)$へ移るので，拡大は非分裂である．
異なるcharacterを持つ二次元非分裂表現の像には，二つの対角冪等元と非零の非対角cornerがあり，
像は全三角行列代数である．

同様に$\Ext_A^1(W,U)=\Ext_X^2(P,S_-)=\Hom_X(S_-,P)^*=k$から$V'$が存在する．
$V\oplus V'$への像は，共通の二つの対角成分と両向きの非対角成分からなる
四次元代数$C_0\simeq\Pi(A_2)$の部分代数$C$である．characterの分離から
$C$の半単純像は$k\times k$である．$J=\operatorname{rad}C_0$は$J^2=0$なので，
冪等元を持ち上げ，$1+J$による共役で標準の対角冪等元に合わせる．
両拡大の非分裂性により二つの一次元cornerがともに像へ入り，$C=C_0$となる．
最後に$[P]=[S_+]+[S_-]$，$[S_-[1]]=-[S_-]$を使う．
\end{proof}

反射$s_\beta(v)=v+(v\beta)\beta$は$(\alpha,\beta)$を$(\alpha+\beta,-\beta)$へ送る．
従ってこの局所模型では，正負の共鳴から反標準heartの正指数二単純への組替えがWeyl反射に対応する．
これは全$E_6$のheartを一つの反射で記述する主張ではない．
商$A\to\Pi(A_2)$もhomological epimorphismではない．商の単純加群の高次Extは二周期的だが，
ambientのExtは次数2を超えて消えるからである．

\begin{remark}[同型型と生成元の係数]
この定理は実際の有限次元商の同型型を与える．しかし，楕円データからあらかじめ固定した
$A$の生成元に対するcharacterの値や拡大cocycleの係数を数値式で与えてはいない．
その明示計算，および全$E_6$の境界外圏のアフィン作用は別の課題として残る．
\end{remark}

\section{四者を同じ図式に置くための条件}\label{sec:alignment}
比較は二圏ずつの同値の存在だけでは完了しない．同じ$e$と$s$を使って反標準環を作り，
比較の合成がその構造を保つことを確認する必要がある．
\begin{proposition}[pointedな比較の接続]\label{prop:alignment}
次のデータが与えられたとする．
同じblow-up表示によるpointedな$J:\Db(\AV)\simeq\Db(\AR)$，
定理\ref{thm:derived}の$\Phi$，およびpointedな
$E:\Db(\AAS)\simeq\Db(\AR)$がある．さらにAS側で系\ref{cor:ASac}の条件と
命題\ref{prop:unique}の高次消滅が成立するとする．
このとき$\Phi E$は標準AS heartから反標準heartへのt-exact同値であり，
$\Phi J$はVdB heartに対して$J^{-1}(\cT),J^{-1}(\cF)$によるHRS比較となる．
従って$\cF=0$なら四者はこれらの指定関手の下で一致する．
\end{proposition}
\begin{proof}
corner同値からAS側の指定反標準orbitは豊富である．pointed性により環とorbitを同定でき，
両heartは命題\ref{prop:unique}の認識式を満たすので$\Phi E$はt-exact．
$J$のt-exactnessはtorsion pairを$\AV$へ運び，定理\ref{thm:globaltilt}の式を保つ．
$\cF=0$のときの結論は命題\ref{prop:hrs}から従う．
\end{proof}

AOUの許容領域では定理\ref{thm:AOU-pointed}によりsection model--AS比較のpointed性が確定した．
これを上の$E$へ接続するには，同じsection modelとRainsデータとの指定した同定を用いる．
確認済みの範囲と残る同定は，次の表のとおりである．
\begin{center}
\small
\begin{tabular}{p{.18\textwidth}p{.36\textwidth}p{.35\textwidth}}
\toprule
比較 & 結論と適用範囲 & 自動的には従わないこと\\
\midrule
VdB--Rains & 同じblow-up・gluing・blow-down表示でabel同値 & 独立に選んだavatarのデータ対応\\
Rains--反標準 & genuine smooth楕円gluingで導来同値，一回のHRS tilt & 任意nc曲面・特異境界への拡張\\
VdB--section model & VdBの一般位置条件下で振幅1の導来同値 & raw blow-upとのheart一致\\
AS--反標準 & foldingと両bimodule有限性の下，有限cornerとabel同値が同値 & 全AS locusでの条件成立\\
section model--AS & 許容領域でpointed導来同値，反標準環も一致 & 全域での標準heartの一致\\
四者の接続 & pointed比較を揃えれば命題\ref{prop:alignment} & 相対familyと標識忘却の降下\\
\bottomrule
\end{tabular}
\end{center}

\subsection{可換基準点と非可換共鳴}
滑らかな可換三次曲面$S$では，実際のblow-up heartとRainsの可換heartは$\coh S$である．
$-K_S$がampleなので反標準環の$\qgr$も$\coh S$を回復する．
標準的なline bundle helixの各周期部分列$\OO(D_i-nH)$もSerreの生成・消滅を満たし，
ASの$\mathbb Z$-代数から同じheartを再構成する．この場合には四者の一致を実際に確認できる．
しかし，一つの可換点での一致から，非可換パラメータ空間での一致locusの開性を推論しない．
命題\ref{prop:density}は，楕円chartでは負指数共鳴が可算個の軌跡の合併として稠密となることを示す．
「稠密な共鳴」と「一つの閉じた判別因子」，あるいは「導来特異化」は異なる概念である．

\subsection{族に移すときの追加課題}
本稿の主要な比較は幾何学的点での結果である．基礎$T$上の族へ移すには，
少なくともgluing bimoduleとkernelの相対perfect性，orbitのHomの基底変更，
有限性と消滅次数の一様性，t構造の基底変更適合性を確認する必要がある．
構造対象の同型と反Serre作用への適合を降下データとして保つことも必要である．
構造対象の同定については命題\ref{prop:AOU-relative-object}が相対Homのperfectnessと
基底変更の下で評価triangleから自然な同型を与える．その前提となる相対圏の構成と
全heartの降下を，この対象単独の同定から推論しない．

有限cornerの範囲だけなら次の初等的な相対判定が使える．
\begin{proposition}\label{prop:relativecorner}
有限型基礎$T$上に次数1生成のgraded代数族$C$があり，各$C_N$が有限表示で
fiberの形成と両立するとする．このときfiberが有限次元となるlocusは
$\bigcup_{N\ge1}\{t:(C_N)_t=0\}$であり開である．
$T$がnoetherianなら，この開locus上である共通の$N$を選べる．
\end{proposition}
\begin{proof}
各集合はcoherent module $C_N$のsupportの補集合であり開である．
補題\ref{lem:cutoff}によってこれらは$N$とともに増加し，合併が求めるlocusとなる．
noetherian空間の開集合は準コンパクトなので有限部分被覆を選び，最大の$N$を取る．
\end{proof}
この命題は有限cornerのfiber条件を述べるだけであり，heartの相対同値・flatness・stack降下を
証明するものではない．本稿のt構造比較と，内在的非可換del Pezzoモジュライの
全成分分類は分離して扱う．

\appendix
\section{三次曲面の数値データと再現手順}\label{app:numerics}
可換三次曲面の標準格子を
$\Pic(S)=\mathbb Zl\oplus\bigoplus_{i=1}^6\mathbb Ze_i$，
$l^2=1$，$e_i^2=-1$，$H=3l-\sum e_i$とする．
AOU \cite[(4.1)]{AOU}の九線束のChern類は
\[
0,\quad -2l+\sum e_i,\quad -l+e_1+e_2+e_3,
\quad e_1,e_2,e_3,\quad l-e_4,l-e_5,l-e_6.
\]
各三つが一blockをなし，$H$次数は$0,1,2$である．
$\chi(\OO(D),\OO(E))=1+\tfrac12(E-D)(E-D+H)$からEuler行列は
\[
C=\begin{pmatrix}I_3&J_3&2J_3\\0&I_3&J_3\\0&0&I_3\end{pmatrix},
\qquad \det C=1,
\]
ただし$J_3$は全成分1の行列．成分総和は45である．
これは可換点の既知のfull strong性と合わせてEnd代数の次元を確認するが，
Euler行列だけからstrong性を証明してはいない．
$K_0$でのSerre作用は$C^{-1}C^{\mathsf t}$，反Serre作用はその逆であり，
後者は$H$によるtensorと一致する．

根は$e_i-e_j$ $(i\ne j)$，$\pm(l-e_i-e_j-e_k)$ $(i<j<k)$，
$\pm(2l-\sum e_i)$の72個である．各根について$\alpha^2=-2$，$H\alpha=0$．
反標準Euler値は$n=0,\ldots,8$で
\[
1,4,10,19,31,46,64,85,109
\]
となる．これらを添付コードで有理数上の厳密計算として確認する．

混合$A_2$では，四直既約のHom表をquiver射の線形方程式から求め，
Euler形式とSerre双対からambient Ext$^1$を検算する．
指数$p,n=1,\ldots,8$，四直既約の重複度$0,\ldots,4$について
40,000通りのcohomology公式を検査する．$p=n$の検算は仮想的線形代数の確認であって，
genuine avatar上にそのsheafが存在することの証拠ではない．

構造対象のmutationについては，AOUの操作の合成順序を実装し，
標識対象の位置が$9,9,9,6,6,3,1$と移ることと，最後の二箇所だけに
Hom消滅の入力が必要であることを検査する．
\eqref{eq:two-degree-zero}の指数ベクトルと次数零性も整数演算で検査する．
これは非自明な楕円線束のcohomology消滅やmutation対象の同型を計算で証明するものではなく，
証明に使う順序と数値の独立な確認である．

本稿は主TeXファイルと四つのsectionソースからなり，\texttt{latexmk -lualatex}で再生成できる．
日本語LuaLaTeX，LuaTeX-ja，\texttt{tikz-cd}，Python~3を用いる．
リポジトリrootでの主要コマンドは
\begin{verbatim}
python3 computations/verify_cubic_heart_numerics.py
python3 computations/verify_four_hearts.py
latexmk -lualatex papers/nc_del_pezzo_hearts.tex
\end{verbatim}
である．数値検算は，sheafの存在，derived type，導来完全忠実性，
または新規性を計算によって証明するものではない．



\begin{thebibliography}{99}
\bibitem{VdB}
M.~Van den Bergh, \emph{Blowing up of non-commutative smooth surfaces},
Mem. Amer. Math. Soc. \textbf{154} (2001), no.~734;
\href{https://arxiv.org/abs/math/9809116}{arXiv:math/9809116v1}.
使用箇所：Lemma 9.1.5，Definition 11.1.2，Theorem 11.1.3，Chapter 12，Remark 12.0.5．
\bibitem{Rains}
E.~M.~Rains, \emph{The birational geometry of noncommutative surfaces},
\href{https://arxiv.org/abs/1907.11301}{arXiv:1907.11301v1} (2019).
使用箇所：Lemma 4.9，Theorem 4.10，Propositions 4.13, 4.17，Theorem 4.30，
\S5，\S7，Theorem 9.10，Corollaries 10.12, 10.16，Proposition 10.13，\S11．
\bibitem{Filtered}
E.~M.~Rains, \emph{Filtered deformations of elliptic algebras},
\href{https://arxiv.org/abs/2107.13540}{arXiv:2107.13540v2} (2022).
使用箇所：Corollary 3.8とtraceの証明，\S4，Proposition 5.1．
\bibitem{RainsBook}
E.~M.~Rains, \emph{The (noncommutative) geometry of difference equations},
\href{https://arxiv.org/abs/2504.16187}{arXiv:2504.16187} (2025).
本稿は書誌と概要を確認し，定理番号・仮定の照合には上記2019年稿を用いた．
\bibitem{AOU}
T.~Abdelgadir, S.~Okawa, K.~Ueda,
\emph{Compact moduli of marked noncommutative cubic surfaces},
\href{https://arxiv.org/abs/2404.00175}{arXiv:2404.00175v1} (2024).
使用箇所：\S3--4，(4.1)，\S9，Lemma 9.2，Theorem 9.4，(9.44)，(9.54)--(9.57)．
\bibitem{OU}
S.~Okawa, K.~Ueda, \emph{Marked Artin--Schelter surfaces of del Pezzo types},
\href{https://arxiv.org/abs/2508.04986}{arXiv:2508.04986v1} (2025).
使用箇所：Definitions 1.1, 3.1，\S6．
\bibitem{AZ}
M.~Artin, J.~J.~Zhang, \emph{Noncommutative projective schemes},
Adv. Math. \textbf{109} (1994), 228--287.
\href{https://doi.org/10.1006/aima.1994.1087}{doi:10.1006/aima.1994.1087}.
\bibitem{CHZ}
X.-W.~Chen, Z.~Han, Y.~Zhou, \emph{Derived equivalences via HRS-tilting},
Adv. Math. \textbf{354} (2019), 106749;
\href{https://arxiv.org/abs/1804.05629}{arXiv:1804.05629v2}.
使用箇所：Theorems A, B，\S2--3．
\bibitem{RSS}
D.~Rogalski, S.~J.~Sierra, J.~T.~Stafford,
\emph{Noncommutative Blowups of Elliptic Algebras},
\href{https://arxiv.org/abs/1308.2216}{arXiv:1308.2216v2} (2015).
使用箇所：Proposition 2.4（正次数，無限位数の境界自己同型）．
\bibitem{YZ}
A.~Yekutieli, J.~J.~Zhang, \emph{Rings with Auslander Dualizing Complexes},
J. Algebra \textbf{213} (1999), 1--51;
\href{https://arxiv.org/abs/math/9804005}{arXiv:math/9804005v3}.
使用箇所：Theorems 4.5, 5.1(1), 5.13．
\bibitem{Orlov}
D.~Orlov, \emph{Smooth and proper noncommutative schemes and gluing of DG categories},
\href{https://arxiv.org/abs/1402.7364}{arXiv:1402.7364v5} (2015).
使用箇所：Theorems 3.18, 3.25，Corollary 3.19．
\end{thebibliography}
\end{document}
